\documentclass[11pt]{article}

\usepackage[a4paper,margin=1in]{geometry}
\usepackage{amsmath,amssymb,amsthm}
\usepackage{stmaryrd}      % \llbracket \rrbracket
\usepackage{mathtools}
\usepackage{enumitem}
\usepackage{xcolor}
\usepackage[most]{tcolorbox}
\usepackage{hyperref}
\hypersetup{colorlinks=true,linkcolor=blue!50!black,urlcolor=blue!50!black,
  citecolor=blue!50!black}

% ---- notation ----------------------------------------------------------
\newcommand{\comp}{\mathbin{\triangleleft}}       % composition product of Poly
\newcommand{\Set}{\mathsf{Set}}
\newcommand{\Poly}{\mathsf{Poly}}
\newcommand{\Cat}{\mathsf{Cat}^{\sharp}}
\newcommand{\Comod}{\mathsf{Comod}}
\newcommand{\PRA}{\mathsf{PRA}}
\newcommand{\St}{\mathsf{St}}                      % costate (store) comonoid
\newcommand{\Hom}{\mathrm{Hom}}
\newcommand{\I}{\mathbb{I}}                        % indiscrete category
\newcommand{\y}{\mathsf{y}}
\newcommand{\bk}[1]{{#1}^{\sharp}}                 % backward part of a morphism
\newcommand{\const}{\mathrm{const}}
\newcommand{\NC}{\mathrm{NC}}
\DeclareMathOperator*{\colim}{colim}

% ---- theorem environments ----------------------------------------------
\theoremstyle{plain}
\newtheorem{theorem}{Theorem}
\newtheorem{lemma}{Lemma}
\newtheorem{proposition}{Proposition}
\newtheorem{corollary}{Corollary}
\theoremstyle{definition}
\newtheorem{definition}{Definition}
\newtheorem{example}{Example}
\theoremstyle{remark}
\newtheorem{remark}{Remark}

\newtcolorbox{mainbox}{colback=blue!4,colframe=blue!40!black,
  boxrule=0.6pt,arc=1pt,left=6pt,right=6pt,top=4pt,bottom=4pt}

\newcommand{\contentcommit}{a452538}

\title{\textbf{The worker is the constant-branch slice of the weld}\\[3pt]
  \large A bicomodule comparison of state-entanglement and data-aggregation}
\author{MacBeth\\[-2pt]\small Kodamai}
\date{19 October 2026}

\begin{document}
\maketitle

\begin{center}
\fbox{\begin{minipage}{0.93\textwidth}
\small
\textbf{Author:} MacBeth (Kodamai).\quad\textbf{Date:} 2026-10-19.\\
\textbf{Prepared for} Rick (referee); \textbf{CC} Robin Langer.\quad
\textbf{Content commit:} \texttt{\contentcommit}.\\
\textbf{Status.} Round-3 referee response; supersedes the round-2 content \texttt{a452538}.
Registry node \texttt{weld-is-regular-bicomodule}: root \textbf{[proved]} (clause-(1) core
peer-reviewed); \texttt{clause3} \textbf{[peer-reviewed]} (round-3 endorsement), \texttt{clause2}
\textbf{[peer-reviewed]} on delivery of the round-3 \textsc{must}~M1 fix carried by this revision.
Self-verification of the retract claim (M1): \texttt{scratch/write-2026-10-18-M1-selfcheck.py};
machine-checked in \texttt{lean/2026-10-18-retract-bicomodule-map.lean} (axioms
\texttt{[propext]}). Source: Spivak--Garner--Fairbanks, \emph{Functorial Aggregation},
arXiv:2111.10968 (``SGF'').
\end{minipage}}
\end{center}

\begin{abstract}
In the polynomial-functor semantics of stateful agents, two constructions package the
state of a worker: the \emph{weld} $B=\St S\comp\St T$, which composes two costate
comonads, and the \emph{worker grading} $A=\St(S\times T)$, which tensors them.
It was recently established that $B$ is a genuine $(\St T,\St S)$-bicomodule whose
represented parametric-right-adjoint functor is $\bar F_B(V)=(T\times V)^S$. The status
of the worker $A$ inside this bicomodule was left open, and a natural conjecture --- that
$A$ is the right-adjoint ``content'' factor of $B$ --- had resisted a clean verdict,
because $A$ and $B$ live at different levels of the framed bicategory $\Cat$ and so cannot
be compared as written. We resolve this by reading $A$ in the same variance as $B$: the
worker is canonically the \emph{constant-branch sub-bicomodule} $B'\subseteq B$, with
carrier polynomial exactly $A$. Worker and weld then \emph{share} the content power
$(-)^S$ and differ \emph{only} in their coproduct-aggregation index --- the constant branch
maps $T$ for the worker against all branch maps $T^S$ for the weld. The separation is
proper precisely on $|S|,|T|\ge 2$, and degenerates (worker $=$ weld on the nose) when
$|S|=1$ or $|T|=1$. State-entanglement and data-aggregation are thus genuinely different
collapses of one bicomodule.
\end{abstract}

\section{Introduction}

\paragraph{The problem.}
A stateful agent --- a \emph{worker} --- is modelled in the container/polynomial program
as an interface together with a store of internal state. When a worker with input set
$S$ and output set $T$ carries its state, two polynomial constructions compete to express
``state over the interface''. The first is the \emph{weld}
\[
  B \;=\; \St S \comp \St T,
\]
the composition-product ($\comp$) of the costate comonads of $S$ and $T$; this is the
Zappa--Sz\'ep style orchestration weld, in which $S$ runs and then $T$ runs. The second is
the \emph{worker grading}
\[
  A \;=\; \St S \otimes \St T \;=\; \St(S\times T),
\]
the Dirichlet tensor of the two comonads, which fuses the two states into a single store
on $S\times T$. Both are comonoids --- small categories --- in $(\Poly,\comp,\y)$, and both
claim to be ``the state of the composite worker''. Which is right, and how are they
related?

\paragraph{The gap.}
The question acquired teeth once the weld was placed inside a framed bicategory.
By the Spivak--Garner--Fairbanks correspondence~\cite{SGF}, the framed bicategory $\Cat$
of categories and parametric right adjoints is equivalent to $\Comod(\Poly)$, so a
horizontal $1$-cell $c\to d$ is a \emph{bicomodule}, equivalently a parametric right
adjoint (PRA) functor $c\text{-}\Set\to d\text{-}\Set$. A companion note established three
facts about the weld, of which the first was independently refereed and confirmed:
\begin{enumerate}[label=(\arabic*),leftmargin=2.2em]
\item \emph{(existence, peer-reviewed)} $B=\St S\comp\St T$ carries a genuine
  $(\St T,\St S)$-bicomodule structure, hence is a PRA functor
  $\St T\text{-}\Set\to\St S\text{-}\Set$, with
  \[
     \bar F_B(V)\;\cong\;(T\times V)^S\;=\;T^S\times V^S ;
  \]
\item \emph{(identification)} the worker $A$ is the right-adjoint ``content'' factor of
  $B$ under the SGF slice factorization;
\item \emph{(tail)} the data discarded by the worker's retract is the coproduct
  aggregation tail of $B$.
\end{enumerate}
Clauses (2) and (3) were the live question, and the obstacle was a \emph{type error}:
the worker $A$ is a priori a $0$-cell of $\Cat$ (a comonoid, a category), while the weld
content is a $1$-cell (a bicomodule $\St T\to\St S$). One cannot ask whether a category
``equals'' the content of a bicomodule until one has read the category as a bicomodule of
the same variance. Until that reading is fixed, any comparison of exponents is comparing
objects at different levels. A second, subtler error compounds this one: inferring from the
carrier coincidence $B=A$ at $|S|=1$ that the worker must there be the content factor. Equal
carriers do not entail equal content factors --- at $(|S|,|T|)=(1,m)$ with $m\ge2$ the
carriers coincide, yet the worker is \emph{not} the content factor --- so the coincidence
$B=A$ settles nothing about content.

\paragraph{The result.}
We supply the missing reading and settle (2) and (3) together. The worker has a canonical
$(\St T,\St S)$-bicomodule structure \emph{inside} the weld: it is the sub-bicomodule
carried by the \emph{constant} branch maps.

\begin{mainbox}
\begin{theorem}[Worker $=$ constant-branch slice; see
Proposition~\ref{prop:sigma} and Theorems~\ref{thm:compare}, \ref{thm:retract}]\label{thm:main}
Let $S,T$ be finite nonempty sets. Read the worker $A=\St(S\times T)$ as a
$(\St T,\St S)$-bicomodule via the inclusion $\sigma\colon A\xrightarrow{\ \sim\ }B'$
onto the constant-branch sub-bicomodule $B'\subseteq B$
$(\sigma_1(s,t)=(s,\const_t),\ \bk\sigma=\mathrm{id}_{S\times T})$. Then
\[
  \bar F_A(V)\;=\;\textstyle\sum_{g\in \const(T)} V^S \;=\; T\cdot V^S
  \qquad\subseteq\qquad
  \bar F_B(V)\;=\;\textstyle\sum_{g\in T^S} V^S \;=\; T^S\cdot V^S .
\]
Worker and weld \emph{share} the right-adjoint content power $(-)^S$; they differ
\emph{only} in the coproduct-aggregation index, $T\cong\const(T)$ for the worker against
all branch maps $T^S=\Hom(S,T)$ for the weld. The inclusion $\bar F_A\hookrightarrow\bar
F_B$ is an equality iff $|S|=1$ or $|T|=1$, and is proper precisely on $|S|\ge2\wedge
|T|\ge2$.
\end{theorem}
\end{mainbox}

Thus clause (2) is \textbf{false} where it is interesting ($|S|,|T|\ge2$) --- but
\emph{not} for the reason first proposed. The content \emph{power} $(-)^S$ is in fact
\emph{shared}; the worker is a \emph{sub-coproduct} of the weld's content, $T$ copies against
$T^S$. We are careful to separate two readings that the companion note conflated: the worker
is the bare \emph{content factor} $(-)^S$ only when $|T|=1$, whereas it equals the
\emph{whole} weld (carriers $A=B$) when $|S|=1$ or $|T|=1$ --- so at $(1,2)$ the worker is the
whole weld yet not the content factor. And clause (3), once stated honestly, splits into a
true half (the aggregation tail \emph{is} the full branch-map set $T^S$) and a false half
(the retract does \emph{not} discard ``the non-constant maps''); moreover the set-level
collapse $e$ realising it is not itself a bicomodule map --- the genuine bicomodule retraction
is a base-pointed variant $r_{s_0}$.

\paragraph{The method.}
The one idea is Definition~\ref{def:slice}: inside the weld's position set $S\times T^S$,
the positions with \emph{constant} branch map form a sub-polynomial whose carrier is
exactly $A$, and this slice is coaction-closed because it is \emph{$S$-saturated}
(Lemma~\ref{lem:subbicomod}(b)): $\lambda$ moves the input over all of $S$ while leaving the
branch map untouched, so a set closed under that motion and built from constant branch maps
stays inside the constant slice. Everything else is bookkeeping on the peer-reviewed formula
$\bar F_B(V)=(T\times V)^S$: restrict its coproduct index from $T^S$ to $\const(T)$
(Section~\ref{sec:compare}), and read off what the worker's retract keeps and discards
(Section~\ref{sec:retract}).

\paragraph{Consequences.}
The two constructions carry a genuine semantic distinction that the orchestration calculus
must respect. State-\emph{entanglement} (the worker $A$) records one branch choice per
output token $t\in T$; data-\emph{aggregation} (the weld $B$) records one branch choice
per response policy $g\in T^S$. Both read the whole input state --- the shared content
$(-)^S$ --- but the weld aggregates over all policies and the worker only over the constant
ones. The exact locus where the distinction bites is $|S|,|T|\ge2$: with a single input or
a single output there is nothing to entangle, and $A=B$.

\paragraph{Status and outline.}
This is a referee-response revision: clause (1) and the formula $\bar F_B(V)=(T\times V)^S$
are peer-reviewed input, used but not re-litigated. Section~\ref{sec:prelim} fixes
conventions (bicomodules, their handedness, the PRA correspondence, the costate comonad).
Section~\ref{sec:weld} recalls the clause-(1) core. Section~\ref{sec:slice} proves the new
content: the worker is the constant-branch sub-bicomodule. Section~\ref{sec:compare}
carries out the corrected comparison, Section~\ref{sec:retract} the corrected retract
analysis. Section~\ref{sec:unified} states the unified result and its orchestration
reading. Section~\ref{sec:response} is a point-by-point response to the referee.

\section{Preliminaries}\label{sec:prelim}

We work in the category $\Poly$ of polynomial functors $\Set\to\Set$ and natural
transformations, embedded fully faithfully in $[\Set,\Set]$~\cite{SGF,NiuSpivak}. A
polynomial $p=\sum_{i\in p(1)}\y^{p[i]}$ has a set of \emph{positions} $p(1)$ and, at each
position $i$, a \emph{direction set} $p[i]$; it acts by $p(X)=\sum_i X^{p[i]}$. The
\emph{composition product} is functor composition $(p\comp q)(X)=p(q(X))$ with unit $\y$;
the \emph{Dirichlet product} $\otimes$ has positions $p(1)\times q(1)$, directions
$p[i]\times q[j]$, and unit $\y$. A morphism $\phi\colon p\to q$ is a function
$\phi_1\colon p(1)\to q(1)$ on positions together with a backward function
$\bk\phi\colon q[\phi_1 i]\to p[i]$ on directions, for each $i$.

\paragraph{Comonoids are categories.}
A comonoid in $(\Poly,\comp,\y)$ is exactly a small category~\cite[Prop.~5.2]{SGF}: the
positions are objects, the directions at an object are the morphisms out of it, the counit
selects identities and the comultiplication is composition. We write $\Cat$ for the
framed bicategory whose objects are these polynomial comonoids.

\begin{definition}[Bicomodule; {\cite[Def.~4.2, Rmk~4.3]{SGF}}]\label{def:bicomod}
Let $(c,\varepsilon_c,\delta_c)$ and $(d,\varepsilon_d,\delta_d)$ be comonoids in
$(\Poly,\comp,\y)$. A \emph{$(c,d)$-bicomodule} $m\colon c\to d$ is a carrier $m\in\Poly$
with a left $d$-coaction and a right $c$-coaction
\[
   d\comp m \;\xleftarrow{\ \lambda\ }\; m \;\xrightarrow{\ \rho\ }\; m\comp c
\]
satisfying counit, coassociativity, and compatibility laws. (The variance follows the
Leibniz reading of~\cite[Rmk~4.3]{SGF}: $\lambda$ sends $m$ to ``$d$ after $m$'' and $\rho$
to ``$m$ after $c$'', so $m$ goes from $c$ to $d$.)
\end{definition}

\begin{theorem}[SGF correspondence; {\cite[Prop.~3.9, Thm.~4.28]{SGF}}]\label{thm:sgf}
For comonoids $c,d$, the framed bicategories $\Cat$ and $\Comod(\Poly)$ are equivalent;
in particular
\[
  \Comod(\Poly)(c,d)\;\simeq\;\PRA\bigl(c\text{-}\Set,\ d\text{-}\Set\bigr),
\]
so a $(c,d)$-bicomodule \emph{is} a parametric right adjoint functor
$c\text{-}\Set\to d\text{-}\Set$. For a $c$-variable polynomial $F$, the SGF slice
factorization writes the PRA as a coproduct of representables indexed by $F(1)$, with
right-adjoint content given by the connected-limit part~\cite[Prop.~3.6]{SGF}.
\end{theorem}

\begin{definition}[Costate comonad; {\cite[Ex.~5.6]{SGF}}]\label{def:costate}
For a nonempty set $S$ write
\[
  \St S \;=\; S\cdot\y^{S} \;=\; \textstyle\bigl[\begin{smallmatrix}S\\S\end{smallmatrix}\bigr],
\]
the \emph{costate} (or store) comonad, i.e.\ the comonad induced by the adjunction
$S\times(-)\dashv(-)^S\colon\Set\to\Set$. It has $S$ positions, each with direction set
$S$. As a category it is the \emph{indiscrete category} $\I(S)$ on object set $S$: a unique
morphism between any two objects.
\end{definition}

Throughout, $S,T$ are finite and nonempty. We reserve the symbol $\Delta$ for the diagonal
functor and never for a comonoid; the two costate comonoids are $\St S$ and $\St T$, and
the worker carrier is $\St(S\times T)$.

\begin{lemma}[The indiscrete collapse]\label{lem:collapse}
For $S\ne\emptyset$ the indiscrete category $\I(S)$ is equivalent to the terminal category
$\mathbf 1$; hence $\St S\text{-}\Set\simeq\Set$, and under this equivalence every
copresheaf $X$ has $\lim X\cong X(s_0)\cong\colim X$ for each object $s_0$, and the
representable at any object is the terminal copresheaf.
\end{lemma}

\begin{proof}
In $\I(S)$ there is a unique isomorphism between any two objects, so the inclusion
$\mathbf 1\to\I(S)$ of any object is an equivalence. Equivalent categories have equivalent
copresheaf categories, and limits/colimits transport accordingly.
\end{proof}

We use $\bar F$ for the endofunctor of $\Set$ obtained by transporting a store-bicomodule's
PRA functor across the collapses $\St S\text{-}\Set\simeq\Set$,
$\St T\text{-}\Set\simeq\Set$; right-adjointness and the SGF factorization are
equivalence-invariant.

\section{The weld bicomodule (peer-reviewed core)}\label{sec:weld}

We recall the clause-(1) results, which are input to what follows. The naming departs from
the companion note in one respect: $B=\St S\comp\St T$ is \emph{not} a regular bicomodule
(the regular bicomodule of a comonoid $c$ is $c$ as a $(c,c)$-bicomodule). It is the
\emph{unit-composite} bicomodule --- the composite over the unit $\y$ of the two one-sided
regular comodules --- as Theorem~\ref{thm:prafun} makes explicit.

\begin{lemma}[Unit-composite bicomodule]\label{lem:unitcomp}
For any comonoids $c,d$ in a monoidal category $(\mathcal V,\comp,\y)$, the object $d\comp
c$ is a $(c,d)$-bicomodule with
\[
  \lambda \;=\; \delta_d\comp c \colon d\comp c\to d\comp d\comp c,
  \qquad
  \rho \;=\; d\comp\delta_c \colon d\comp c\to d\comp c\comp c .
\]
\end{lemma}

\begin{proof}
The left laws (counit, coassociativity) are the comonoid laws of $d$ whiskered on the right
by $\comp\,c$; the right laws are the comonoid laws of $c$ whiskered on the left by
$d\comp$. For compatibility, $\lambda$ acts on the first ($d$) factor and $\rho$ on the last
($c$) factor; since $\comp$ is a bifunctor, the two whiskerings commute and both composites
equal $\delta_d\comp\delta_c$.
\end{proof}

\begin{theorem}[Clause (1), peer-reviewed]\label{thm:weld}
The weld $B=\St S\comp\St T$ is the unit-composite $(\St T,\St S)$-bicomodule (take
$c=\St T$, $d=\St S$, carrier $d\comp c=\St S\comp\St T$), with $\lambda=\delta_{\St
S}\comp\St T$ and $\rho=\St S\comp\delta_{\St T}$. Hence $B$ is a horizontal $1$-cell of
$\Cat$, i.e.\ a PRA functor $\St T\text{-}\Set\to\St S\text{-}\Set$.
\end{theorem}

\begin{proof}
Lemma~\ref{lem:unitcomp} with $c=\St T$, $d=\St S$; both are comonoids, and
Theorem~\ref{thm:sgf} supplies the PRA reading.
\end{proof}

\begin{theorem}[The represented prafunctor; peer-reviewed]\label{thm:prafun}
After the collapse of Lemma~\ref{lem:collapse},
\[
  \bar F_B(V)\;\cong\;(T\times V)^S\;=\;T^S\times V^S .
\]
As a polynomial endofunctor of $\Set$, $\bar F_B=T^S\cdot\y^S$: its position set (the
coproduct index $\bar F_B(1)$) is the branch-map set $T^S=\Hom(S,T)$, and its direction set
is $S$.
\end{theorem}

\begin{proof}
The unit-composite factors over $\y$ as
\[
   \St T \xrightarrow{\ G\ } \y \xrightarrow{\ H\ } \St S,
   \qquad H\comp G \;=\; \St S\comp\St T \;=\; B,
\]
where $G$ is the right $\St T$-comodule with carrier $\St T$ and $\rho=\delta_{\St T}$, and
$H$ is the left $\St S$-comodule with carrier $\St S$ and $\lambda=\delta_{\St S}$.
Bicomodule composition over $\y$ is just $\comp$, so $F_B=F_H\circ F_G$.

For $F_G\colon\St T\text{-}\Set\to\Set$, the right $\St T$-comodule with carrier $\St T$ is
\[
   F_G(X)\;=\;\sum_{t\in T}\St T\text{-}\Set(\y_t,X),
\]
where $\y_t$ denotes the representable at the object $t$. By Lemma~\ref{lem:collapse} each
$\y_t$ is terminal, so the hom-set is $\lim X$; hence $F_G(X)=\sum_{t\in T}\lim X$ and,
after transport, $\bar G(V)=T\times V$.

For $F_H\colon\Set\to\St S\text{-}\Set$, a left $\St S$-comodule is a diagram
$\St S\to\Poly$~\cite[Prop.~5.7]{SGF}; the regular one assigns to each object $a\in S$ the
component polynomial $\y^S$ (one position, $S$ directions, read off $\delta_{\St S}$). So
$F_H(W)(a)=(\y^S)(W)=W^S$, and after transport $\bar H(W)=W^S$.

Therefore $\bar F_B(V)=\bar H(\bar G(V))=(T\times V)^S$, and $(T\times V)^S=\{S\to T\times
V\}=\{(g,h):g\colon S\to T,\ h\colon S\to V\}=T^S\times V^S$.
\end{proof}

By Theorem~\ref{thm:sgf} the SGF slice factorization of $\bar F_B(V)=\sum_{g\in T^S}V^S$
has coproduct (``$\Sigma$'') index $\bar F_B(1)=T^S$ and right-adjoint \emph{content} the
power $(-)^S$ (represented by the direction set $S$). This content split is the backbone of
the comparison below.

\begin{remark}
The represented prafunctor $\bar F_B=T^S\y^S$ is neither the underlying polynomial of
$A=(S\times T)\y^{S\times T}$ nor that of the weld carrier $B=(S\cdot T^S)\y^{S\times T}$.
The PRA functor of a bicomodule is not the underlying functor of its carrier --- the trap
at the heart of the problem. (At $|S|=|T|=2$, $V=2$: $\bar F_B=16$, $A_{\mathrm{und}}=64$,
$B_{\mathrm{und}}=128$.)
\end{remark}

\section{The worker as a bicomodule: the constant-branch slice}\label{sec:slice}

Here is the new content. To compare the worker $A$ with the weld content we must read $A$
as a $(\St T,\St S)$-bicomodule of the \emph{same} variance as $B$. There is a canonical
such reading, and it sits visibly inside $B$.

Recall the weld carrier as a polynomial:
$(\St S\comp\St T)(X)=S\times(T\times X^T)^S=S\times T^S\times X^{S\times T}$, so $B$ has
position set $S\times T^S$ --- a current input $s\in S$ and a branch map $g\colon S\to T$
--- and direction set $S\times T$ at each position.

\begin{definition}[Constant-branch slice]\label{def:slice}
Let $\const\colon T\hookrightarrow T^S$, $t\mapsto\const_t$, be the inclusion of the
constant branch maps. Let $B'\subseteq B$ be the sub-polynomial on the positions with
constant branch map,
\[
  B' \;=\; \bigl\{(s,g)\ :\ g\in\const(T)\bigr\},
\]
with inherited direction set $S\times T$. Its position set is $S\times\const(T)\cong S\times
T$, so as a polynomial
\[
  B' \;=\; (S\times T)\cdot\y^{\,S\times T}\;=\;\St(S\times T)\;=\;A
\]
on the nose.
\end{definition}

\begin{lemma}[Branch decomposition and the closed subobjects]\label{lem:subbicomod}
For each branch map $g\in T^S$ let $B_g\subseteq B$ be the positions with that branch map,
$B_g=\{(s,g):s\in S\}$, a sub-polynomial with $|S|$ positions and directions $S\times T$.
Then:
\begin{enumerate}[label=\textup{(\alph*)},leftmargin=2em]
\item each $B_g$ is a $(\St T,\St S)$-sub-bicomodule of $B$, and
  $B=\coprod_{g\in T^S}B_g$ as bicomodules; all the $B_g$ are mutually isomorphic;
\item a sub-polynomial of $B$ is coaction-closed iff it is \emph{$S$-saturated}, i.e.\ a
  union $\coprod_{g\in G}B_g$ for some $G\subseteq T^S$;
\item $B'=\coprod_{g\in\const(T)}B_g$ is coaction-closed; its carrier is $A$.
\end{enumerate}
\end{lemma}

\begin{proof}
The two coactions act on the position $(s,g)$ as follows. The left coaction
$\lambda=\delta_{\St S}\comp\St T$ \emph{moves the input} $s$ over all of $S$ while leaving
the branch map $g$ untouched: $\delta_{\St S}$ replicates the position across the input set,
$\lambda(s,g)=\bigl(s,\ a\mapsto(a,g)\bigr)$, the inner copies ranging over $a\in S$ with
the \emph{same} $g$. The right coaction $\rho=\St S\comp\delta_{\St T}$ acts on the outer
$\St S$-copy and reorganizes the inner $\St T$-layer by $\delta_{\St T}$; it fixes the whole
position $(s,g)$, in particular $g$. (This is the corrected reading: $g$ never appears as an
argument of $\delta_{\St T}$, and it is $\lambda$, not $\rho$, that moves $s$.)

Hence a position set is closed under both coactions iff, whenever it contains $(s,g)$, it
contains $(a,g)$ for every $a\in S$ --- $S$-saturation --- which is exactly (b); and each
$B_g$ (a single $g$, all inputs) is a minimal closed piece, giving the decomposition (a).
The $B_g$ are mutually isomorphic by relabelling $g$. For (c), $\const(T)\subseteq T^S$
yields the $S$-saturated union $B'=\coprod_{t\in T}B_{\const_t}$; it has $|T|\cdot|S|$
positions with directions $S\times T$, i.e.\ carrier $(S\times T)\cdot\y^{S\times T}=A$. The
counit and (co)associativity laws restrict because they are the comonoid laws of
$\St S,\St T$ whiskered, and $B'$ is a full sub-on-positions with inherited directions.
\end{proof}

\begin{remark}[what singles out $B'$]\label{rem:yoneda}
Closure does \emph{not} single out $B'$: by Lemma~\ref{lem:subbicomod}(b) every subset
$G\subseteq T^S$ gives a closed sub-bicomodule, and there are $2^{|T^S|}$ of them. What makes
$\const(T)$ canonical is \emph{naturality}. The constant assignment is the diagonal
$\Delta\colon T\to T^S=\Hom(S,T)$, $t\mapsto\const_t$, the unique map with
$\mathrm{ev}_s\circ\Delta=\mathrm{id}_T$ for every input $s\in S$ --- the unique natural
(Yoneda) section of every evaluation. So among all $S$-saturated slices, $B'$ is the one
indexed by the canonical diagonal $T\hookrightarrow T^S$, and it is this that earns $A$ its
reading, not the fact (shared by many slices) of being closed.
\end{remark}

\begin{proposition}[The worker's canonical reading, via cobase change]\label{prop:sigma}
Equip the worker carrier $A=\St(S\times T)$ with the $(\St T,\St S)$-bicomodule structure
$\pi_!A$ obtained by \emph{cobase change} of the regular $(A,A)$-bicomodule along the two
comonoid projections $\pi_S\colon A\to\St S$ and $\pi_T\colon A\to\St T$ induced by the set
projections of $S\times T$ \textup{(}corestriction of the left coaction along $\pi_S$, of
the right coaction along $\pi_T$\textup{)}. Then the carrier identity
$\sigma\colon A\to B$, $\sigma_1(s,t)=(s,\const_t)$, $\bk\sigma=\mathrm{id}_{S\times T}$, is
an isomorphism of $(\St T,\St S)$-bicomodules $\pi_!A\xrightarrow{\ \sim\ }B'$.
\end{proposition}

\begin{proof}
Cobase change produces a genuine $(\St T,\St S)$-bicomodule on the carrier $A$ with no
reference to $B$, so the comparison is not circular: the left coaction is
$A\xrightarrow{\delta_A}A\comp A\xrightarrow{\pi_S\comp A}\St S\comp A$ and the right is
$A\xrightarrow{\delta_A}A\comp A\xrightarrow{A\comp\pi_T}A\comp\St T$. The two comonoid
projections, as cofunctors, have direction parts $\bk{\pi_S}(a)=(a,t)$ and
$\bk{\pi_T}(t')=(s,t')$ (a direction of the codomain is carried back to one of the domain,
keeping the untouched coordinate fixed). Computing $\lambda_A=(\pi_S\comp A)\,\delta_A$ and
$\rho_A=(A\comp\pi_T)\,\delta_A$ on positions and directions gives, in closed form for
\emph{all} $S,T$,
\[
  \begin{aligned}
    \lambda_A(s,t)&=\bigl(s,\ a\mapsto(a,t)\bigr), & \bk{\lambda_A}\bigl(a,(a',t')\bigr)&=(a',t'),\\
    \rho_A(s,t)&=\bigl((s,t),\ (a,t')\mapsto t'\bigr), & \bk{\rho_A}\bigl((a,t'),t''\bigr)&=(a,t'').
  \end{aligned}
\]
Now apply $\sigma$. On positions $\sigma_1$ is the bijection $S\times T\cong S\times\const(T)$
onto the $B'$-positions and $\bk\sigma=\mathrm{id}_{S\times T}$ matches directions, so pushing
$\lambda_A,\rho_A$ across $\sigma$ reproduces exactly the restricted weld coactions of
Lemma~\ref{lem:subbicomod} at $g=\const_t$: indeed $\lambda_B(s,\const_t)=\bigl(s,\ a\mapsto
(a,\const_t)\bigr)=\bigl(s,\ a\mapsto\sigma_1(a,t)\bigr)$, and $\rho_B$ fixes the whole
position just as $\rho_A$ does. Hence $\sigma$ intertwines both coactions for all $S,T$, and a
carrier bijection with $\bk\sigma=\mathrm{id}$ intertwining $\lambda,\rho$ is an isomorphism
of bicomodules onto its image $B'$. (The five bicomodule laws for $\pi_!A$ and both
equivariance squares of $\sigma$ were additionally machine-checked at all nine shapes
$n,m\in\{1,2,3\}$, and $\bar F_A(V)=T\cdot V^S$ recomputed from $A$'s own cobase-change
structure, independently of $\sigma$.)
\end{proof}

\section{The comparison: a coproduct-index separation}\label{sec:compare}

\begin{theorem}[Corrected comparison]\label{thm:compare}
Read $A$ and $B$ as $(\St T,\St S)$-bicomodules via Section~\ref{sec:slice}
$(A\cong B'\subseteq B)$. Then
\[
  \bar F_B(V)\;=\;\sum_{g\in T^S} V^S \;=\; T^S\cdot V^S,
  \qquad
  \bar F_A(V)\;=\;\sum_{g\in\const(T)} V^S \;=\; T\cdot V^S .
\]
Both have the \emph{same} content power $(-)^S$ (direction set $S$); they differ
\emph{only} in the coproduct index, $T^S$ for the weld against $T\cong\const(T)$ for the
worker. Consequently $\bar F_A$ is a sub-prafunctor of $\bar F_B$, with
$\bar F_A\cong\bar F_B$ \emph{iff $|S|=1$ or $|T|=1$}, and a \emph{proper} inclusion
exactly on $|S|\ge2\wedge|T|\ge2$.
\end{theorem}

\begin{proof}
$\bar F_B=\sum_{g\in T^S}V^S$ is Theorem~\ref{thm:prafun}: the coproduct index is the weld
positions modulo the left $\St S$-collapse (the $S$-head is killed by indiscreteness),
namely the branch maps $T^S$, and the surviving direction is $S$ (the $T$-direction is
killed by the right $\St T$-collapse). By Proposition~\ref{prop:sigma}, $A\cong B'$ is $B$
restricted on positions to $\const(T)\subseteq T^S$ with directions unchanged; the same
recipe restricts the coproduct index to $\const(T)$ and leaves the content power $(-)^S$
untouched, giving $\bar F_A=\sum_{t\in T}V^S=T\cdot V^S$.

Equality of the two functors for all $V$ is equality of polynomials $T^S\y^S$ and $T\y^S$,
i.e.\ $|T^S|=|T|$, i.e.\ $m^n=m$ where $n=|S|,m=|T|$. For $n,m\ge1$ this holds iff $n=1$ or
$m=1$. Otherwise $T\subsetneq T^S$ and the inclusion $\bar F_A\hookrightarrow\bar F_B$ is
proper.
\end{proof}

\begin{remark}[what this corrects]
The exponents \emph{agree} --- both content powers are $(-)^S$. An earlier argument found
``$S$ versus $S\times T$'', but that compared $B$ under its right $\St T$-coaction against
$A$ read with \emph{no} right coaction (the $\y$-reading). With the coaction present, the
$T$-directions of $A$ collapse into the exponent $S$, not $S\times T$. The consistent
separation is the coproduct index $T$ versus $T^S$, nonzero precisely where non-constant
branch maps exist, i.e.\ on $|S|\ge2\wedge|T|\ge2$. In particular at $|S|=1$ one has
$T^S=T$, so $\const(T)=T^S$ and $\bar F_A=\bar F_B=T\cdot V$; moreover the carriers coincide,
$B=\St S\comp\St T=\y\comp\St T=\St T=\St(1\times T)=A$. There the worker is the whole
weld; it is the content factor only if also $|T|=1$.
\end{remark}

\noindent\textbf{Verdict for clause (2).} The original conjecture asks whether the worker is
the right-adjoint \emph{content factor} of the weld, i.e.\ whether $\bar F_A$ equals the
content power $(-)^S$. We must not conflate this with the worker being the \emph{whole} weld.
Two distinct equalities are in play, and they have different loci:
\begin{itemize}[leftmargin=1.5em]
\item \textbf{worker $=$ content factor:} $\bar F_A(V)=V^S$ (the bare content power), which
  holds \emph{iff $|T|=1$} --- a single output collapses the coproduct index $T$ to a point;
\item \textbf{worker $=$ whole weld:} $\bar F_A=\bar F_B$ as prafunctors, equivalently the
  carriers coincide ($A=B$), which holds \emph{iff $|S|=1$ or $|T|=1$}.
\end{itemize}
These are genuinely different: at $(|S|,|T|)=(1,2)$ the carriers coincide, $A=B$ on the nose,
yet $\bar F_A(V)=2V\ne V=$ content, so the worker is the whole weld but \emph{not} the
content factor. Hence the conjecture ``worker $=$ content factor'' is \textbf{true iff
$|T|=1$} and \textbf{false whenever $|T|\ge2$}; on the interesting locus $|S|,|T|\ge2$ the
worker is neither the content factor nor the whole weld. The content power $(-)^S$ is in all
cases \emph{shared}; what the worker is, is a sub-coproduct ($T$ copies of the content against
the weld's $T^S$).

\begin{example}
At $(|S|,|T|)=(2,2)$, $V=2$: $\bar F_A=2\cdot2^2=8$ and $\bar F_B=(2\cdot2)^2=16$ (proper).
At $(1,2)$: $\bar F_A=\bar F_B=4$; at $(1,3)$: $\bar F_A=\bar F_B=6$ (equal). These are
exactly the shapes where the criterion for worker $\ne$ content factor (``$|T|\ge2$'') and
the criterion for worker $\ne$ weld (``$|S|\ge2\wedge|T|\ge2$'') disagree.
\end{example}

\section{The retract: what the worker forgets}\label{sec:retract}

How does the weld project onto the worker slice $B'$? Here one must separate two levels, and
the companion note conflated them. At the level of \emph{position-sets} there is a clean
idempotent collapse; at the level of \emph{bicomodules} that collapse is not a morphism, and
the genuine retraction is a different, base-pointed map. We record both honestly. Throughout,
$n=|S|,\ m=|T|$, and the weld's coproduct index splits as
\[
  T^S \;=\; \const(T)\ \sqcup\ \NC,
  \qquad |\const(T)|=m,\quad |\NC|=m^{n}-m,
\]
with $\NC$ the non-constant maps.

\paragraph{The set-level idempotent $e$.}
Let $e\colon S\times T^S\to S\times T^S$, $e(s,g)=(s,\const_{g(s)})$: collapse each branch
map to the constant at its value on the \emph{current} input $s$.

\begin{theorem}[Keep/discard, as a map of position-sets]\label{thm:retract}
$e$ is idempotent with image $S\times\const(T)\cong S\times T$, the worker positions
$\sigma(A)$. For each current input $s$, $e$ \emph{keeps} the value $g(s)\in T$ (the
diagonal $T$-coordinate) and \emph{discards} the restriction $g|_{S\setminus\{s\}}\in
T^{S\setminus\{s\}}$, a fibre of size $m^{n-1}$. Concretely the fibre of $e$ over
$(s,\const_t)$ is $\{g:g(s)=t\}\cong T^{S\setminus\{s\}}$.
\end{theorem}

\begin{proof}
$e\circ e=e$ since $\const_{g(s)}$ evaluated at $s$ is $g(s)$. The image is the constant
slice $\{(s,\const_t)\}$. Two maps $g,g'$ satisfy $e(s,g)=e(s,g')$ iff $g(s)=g'(s)$; fixing
$g(s)=t$ leaves $S\setminus\{s\}$ free, so the fibre is $T^{S\setminus\{s\}}$, of size
$m^{n-1}$.
\end{proof}

\paragraph{But $e$ is not a bicomodule map.}
Theorem~\ref{thm:retract} is purely combinatorial; it does \emph{not} make $e$ a morphism of
$(\St T,\St S)$-bicomodules, and in fact $e$ is not one.

\begin{proposition}[$e$ fails $\lambda$-equivariance]\label{prop:efail}
With $\bk e=\mathrm{id}$ on directions, $e$ satisfies the right ($\rho$-)equivariance square
on all of $B$, but fails the left ($\lambda$-)equivariance square
$\lambda\circ e=(\St S\comp e)\circ\lambda$ at every position with a non-constant branch map;
hence $e$ is a bicomodule morphism iff every $g\in T^S$ is constant, i.e.\ iff $n=1$ or
$m=1$. The first failure is at $(n,m)=(2,2)$, position $(s,g)=(0,(0,1))$.
\end{proposition}

\begin{proof}[Reason]
$\lambda=\delta_{\St S}\comp\St T$ \emph{moves the input} over all of $S$
(Lemma~\ref{lem:subbicomod}). The map $e$ reads $g$ at the \emph{current} input $s$; after
$\lambda$ has moved the input to some $a\ne s$, the composite $(\St S\comp e)\circ\lambda$
broadcasts $g(a)$, whereas $e$ applied first broadcasts $g(s)$ and $\lambda$ merely copies
it. On the witness $g=(0,1)$ the two give branch value $g(1)=1$ against $g(0)=0$ in the
moved copy --- unequal, so the square fails. ($\rho$ fixes the whole position, so it never
distinguishes the two orders.) A $\mathtt{decide}$-checked verification over the whole carrier
(ground set of size $2$) confirms failure at $(2,2),(2,3),(3,2)$, and confirms the square
holds on the boundary shapes $(1,2),(2,1)$ and on the slice $B'$.
\end{proof}

\paragraph{The genuine bicomodule retraction $r_{s_0}$.}
The fix is to read $g$ at a \emph{fixed} input, immune to $\lambda$'s motion. Choose a base
point $s_0\in S$ and set $r_{s_0}(s,g)=(s,\const_{g(s_0)})$, again with $\bk r=\mathrm{id}$.

\begin{theorem}[The retraction $B\twoheadrightarrow B'$]\label{thm:retract-r}
For any base point $s_0\in S$, $r_{s_0}$ is a morphism of $(\St T,\St S)$-bicomodules
$B\to B'$ \textup{(}it satisfies both the $\lambda$- and $\rho$-equivariance squares\textup{)},
is idempotent, has image the constant slice $B'=\sigma(A)$, and restricts to the identity on
$B'$ --- a genuine bicomodule retraction. Its fibre over $(s,\const_t)$ is
$\{g:g(s_0)=t\}\cong T^{S\setminus\{s_0\}}$, of size $m^{n-1}$: the \emph{same count} as $e$.
Being a bicomodule map, $r_{s_0}$ represents a genuine endomorphism of $\bar F$; the set-level
$e$ does not (Proposition~\ref{prop:efail}), but $e$ agrees with $r_{s_0}$ on the fibre of
$\bar F$ at the base object $s_0$, the point at which the collapse $\St S\text{-}\Set\simeq\Set$
is evaluated. The retraction $r_{s_0}$ is honestly \emph{non-canonical}: it depends on the
choice of $s_0$.
\end{theorem}

\begin{proof}
Because $r_{s_0}$ reads $g$ at the fixed $s_0$, both $r_{s_0}\circ\lambda$ and
$\lambda\circ r_{s_0}$ broadcast $g(s_0)$, so the $\lambda$-square holds; $\rho$ fixes the
position and the $\rho$-square holds as for $e$. Idempotency, image, retraction, and the
fibre count $m^{n-1}$ are as in Theorem~\ref{thm:retract} with $s$ replaced by $s_0$. All of
$\lambda$- and $\rho$-equivariance, idempotency, image$\,=\!B'$, and
$r_{s_0}\circ\iota_{B'}=\mathrm{id}$ were $\mathtt{decide}$-checked at
$(2,2),(2,3),(3,2),(1,2),(2,1)$ (and base-point independence at $s_0\in\{0,1\}$ for $(2,2)$).
\end{proof}

\noindent So the bicomodule-level statement ``the worker is the retract of the weld'' is true
with $r_{s_0}$ in place of $e$; the set-level combinatorics (Theorem~\ref{thm:retract}) are
unaffected, and $e$ agrees with $r_{s_0}$ on the $\bar F$-fibre at $s_0$ (though only
$r_{s_0}$, being a bicomodule map, represents an endomorphism of $\bar F$).

\paragraph{Re-decision of clause (3).}
As first stated --- ``the discarded data is precisely the non-constant branch maps, and that
is the $\Sigma$-tail'' --- the claim is \textbf{false}, for two independent reasons.
\begin{enumerate}[label=(\roman*),leftmargin=2em]
\item The $\Sigma$-tail index is \emph{all} of $T^S=\const(T)\sqcup\NC$, not $\NC$.
  ``$T^S=$ the non-constant maps'' is a cardinality contradiction ($m^n\ne m^n-m$).
\item What the retraction forgets per position is the fibre $T^{S\setminus\{s_0\}}$ of size
  $m^{n-1}$ --- neither $\NC$ (size $m^n-m$) nor $T^S$ (size $m^n$). At $(2,3)$: $3$ versus
  $6$ versus $9$, all distinct.
\end{enumerate}

\noindent The \textbf{true residue} is this: the weld's $\Sigma$-tail \emph{is} the full
branch-map set $T^S=\Hom(S,T)$ --- that half of clause (3) stands --- and the retraction
$r_{s_0}$ realises the projection of this index onto the worker's constant slice
$\const(T)\cong T$ by ``read the branch at $s_0$, then make constant''. The genuinely
discarded structure is the off-diagonal dependence $g|_{S\setminus\{s_0\}}$. It is
\emph{non-trivial} --- the retraction actually loses something, the fibre having size
$m^{n-1}>1$ --- exactly on $|S|\ge2\wedge|T|\ge2$, the same locus as the clause-(2)
separation.

This also settles the old worry about a ``reversal of adjointness'' honestly: $T^S$ is the
index of the \emph{left}-adjoint coproduct $\sum_{T^S}\dashv\Delta_{T^S}$, not right-adjoint
content; the right-adjoint content $(-)^S$ is \emph{kept} by both worker and weld; and the
worker-versus-weld difference is purely how much of the coproduct index is retained ($T$
versus $T^S$) --- no adjoint-direction clash at all.

\section{The unified statement}\label{sec:unified}

\begin{theorem}[Worker $=$ constant-branch slice of the weld]\label{thm:unified}
Let $S,T$ be finite nonempty. With $A$ read as the $(\St T,\St S)$-bicomodule
$\sigma(A)=B'\subseteq B$ of constant branch maps $(\bk\sigma=\mathrm{id}_{S\times T})$,
\[
  \bar F_A(V)=T\cdot V^S \ \subseteq\ \bar F_B(V)=T^S\cdot V^S
\]
is the sub-coproduct over $\const(T)\subseteq T^S$. Worker and weld \emph{share} the
right-adjoint content power $(-)^S$ and differ \emph{only} in the coproduct index ($T$
versus $T^S$). For any base point $s_0$, the bicomodule retraction
$r_{s_0}(s,g)=(s,\const_{g(s_0)})$ projects $T^S$ onto $\const(T)$, forgetting the
off-diagonal fibre $T^{S\setminus\{s_0\}}$ (size $m^{n-1}$); the set-level idempotent
$e(s,g)=(s,\const_{g(s)})$ has the same image and fibre count but is not a bicomodule map.
The degeneracies $\bar F_A=\bar F_B$, trivial discard, and $A=B$ coincide and hold \emph{iff
$|S|=1$ or $|T|=1$}; the proper separation is exactly $|S|\ge2\wedge|T|\ge2$. (The worker is
the bare content factor $(-)^S$ on the narrower locus $|T|=1$.)
\end{theorem}

\paragraph{Orchestration reading.}
State-\emph{entanglement} (the worker $A$) keeps a branch choice per output token $t\in T$;
data-\emph{aggregation} (the weld $B$) keeps a branch choice per response policy $g\in
T^S$. Both read the whole input state --- the shared content $(-)^S$ --- but the weld
aggregates over all policies and the worker only over constant ones. Keeping the two
distinct is a semantic requirement, not a defect, and the exact locus where it matters is
$|S|,|T|\ge2$: with a single input or output there is nothing to entangle.

\paragraph{Grant relevance.}
The weld now has a precise home in $\Poly$: it is a data-migration prafunctor
$\St T\text{-}\Set\to\St S\text{-}\Set$, $X\mapsto(T\times X)^S$, with the branch policy
aggregated by $\sum_{T^S}$. This places the \emph{state} axis of the
compositional-correctness program inside the same framed bicategory $\Cat$ that houses the
directed (retrofunctor) axis. But the worker is not that prafunctor's content: entanglement
and aggregation are different collapses of one bicomodule, and the orchestration semantics
must keep them apart.

\section{Response to the referee}\label{sec:response}

This revision answers the referee's second- and third-round reports. The core finding --- that
$B'$ is coaction-closed and is the worker's reading --- was endorsed; the second round's two
\textsc{must} items and three \textsc{should}s are addressed below, and the third round's
residual items are folded into the same text. We record where each lands.

\begin{itemize}[leftmargin=1.5em]
\item \textbf{M1} (the retract $e$ is not a bicomodule map): Section~\ref{sec:retract}.
  Proposition~\ref{prop:efail} states that $e(s,g)=(s,\const_{g(s)})$ fails
  $\lambda$-equivariance at every non-constant $g$ (first failure $(2,2)$, position
  $(0,(0,1))$), because $e$ reads $g$ at the \emph{current} input while $\lambda$ moves the
  input over $S$. Theorem~\ref{thm:retract-r} replaces it with the base-pointed
  $r_{s_0}(s,g)=(s,\const_{g(s_0)})$, a genuine bicomodule retraction
  $B\twoheadrightarrow B'$ (both squares), honestly non-canonical, keeping every count
  ($m^{n-1}$ fibre) and the same $\bar F$. The set-level facts of Theorem~\ref{thm:retract}
  about $e$ are unchanged; only the bicomodule-retraction reading switches $e\to r_{s_0}$.
  Both the failure of $e$ and the retraction property of $r_{s_0}$ were re-derived by the
  author (companion script; $\mathtt{decide}$-checked in the Lean leaf).
\item \textbf{M2} (clause-(2) verdict conflation): the verdict in
  Section~\ref{sec:compare} now separates the two readings and labels each. ``Worker $=$
  content factor $(-)^S$'' holds \emph{iff $|T|=1$}; ``worker $=$ whole weld'' ($A=B$, equiv.\
  $\bar F_A=\bar F_B$) holds \emph{iff $|S|=1$ or $|T|=1$}; the counterexample $(1,2)$
  ($A=B$ yet $\bar F_A=2V\ne V$) is spelled out. The phrase ``content (indeed $A=B$)'' is
  gone.
\item \textbf{S1} (general sub-bicomodule fact; canonicity): Lemma~\ref{lem:subbicomod} now
  gives $B=\coprod_{g\in T^S}B_g$ with all $B_g$ mutually isomorphic, identifies the
  coaction-closed subobjects as the \emph{$S$-saturated} unions, and
  Remark~\ref{rem:yoneda} singles out $\const(T)$ by the diagonal $\Delta\colon T\to T^S$
  (Yoneda), not by closure --- correcting the earlier ``closed because $g$ is fixed
  pointwise'' mechanism to $S$-saturation.
\item \textbf{S2} ($\sigma$ circular): Proposition~\ref{prop:sigma} now defines the worker's
  structure by \emph{cobase change} $\pi_!A$ of the regular $(A,A)$-bicomodule along
  $\pi_S,\pi_T$ --- a structure built without reference to $B$ --- and exhibits $\sigma$ as
  the carrier isomorphism $\pi_!A\cong B'$. The explicit intrinsic coactions
  $\lambda_A,\rho_A$ are now displayed in the proof of Proposition~\ref{prop:sigma}, so the
  identity $\pi_!A\cong B'$ is established for all $S,T$ (and machine-checked at the nine
  shapes $n,m\in\{1,2,3\}$); no residual coaction-by-coaction item remains.
\item \textbf{S3} (Lemma~3 proof text): the proof of Lemma~\ref{lem:subbicomod} now reads
  ``$\lambda$ moves the input $s$ over $S$ fixing $g$; $\rho$ fixes the whole position'', with
  no spurious $\psi$ in $\delta_{\St T}$.
\item \textbf{N1/N2}: the discarded fibre is called \emph{non-trivial} (size $m^{n-1}>1$),
  not nonempty; the redundant third refutation of clause~(3) has been removed.
\end{itemize}

\noindent The third-round report (a near-endorsement: clause~(3) accepted unconditionally,
clause~(2) accepted once the stray conflation was cleared) raised one \textsc{must} and three
minor items, all textual. Its \textsc{must}~M1 flagged three sentences --- the end of the
``what this corrects'' remark, one sentence of the Introduction, and the last sentence of the
Example of Section~\ref{sec:compare} --- that still ran together ``worker $=$ whole weld''
($A=B$) and ``worker $=$ content factor'' ($\bar F_A=(-)^S$); all three now state the correct
reading, labelling which equality each locus criterion governs ($|T|=1$ for the content
factor, $|S|=1\vee|T|=1$ for the whole weld). Its \textsc{should}~S1 asked that the one
remaining hedge in Proposition~\ref{prop:sigma}'s identification be closed: the explicit
intrinsic coactions $\lambda_A,\rho_A$ are now displayed in that proof and $\pi_!A\cong B'$ is
established for all $S,T$. Its two nits are addressed where they arise: the claim that $e$ and
$r_{s_0}$ ``induce the same map on $\bar F$'' is restricted to the $\bar F$-fibre at the base
object $s_0$ (only $r_{s_0}$, a bicomodule map, represents an endomorphism of $\bar F$), and
the method paragraph now attributes the slice's closure to $S$-saturation
(Lemma~\ref{lem:subbicomod}(b)) rather than to the retired constancy-preservation phrasing.

\noindent Carried over from the first round and unchanged: $B=\St S\comp\St T$ is the
\emph{unit-composite} bicomodule (Lemma~\ref{lem:unitcomp}), not a regular one; the reference
\cite{SGF} is Spivak--Garner--Fairbanks, \emph{Functorial Aggregation} (arXiv:2111.10968),
\emph{not} Niu--Spivak \cite{NiuSpivak}; $S,T$ are standingly nonempty; $\Delta$ denotes the
diagonal, the costate comonoids being $\St S,\St T$.

\paragraph{Honest scope.}
Lemma~\ref{lem:unitcomp} is standard; the contribution is Section~\ref{sec:slice} (the
worker is the constant-branch sub-bicomodule, with carrier exactly $A$, canonical by the
diagonal) together with the corrected clause-(2)/(3) analysis built on the peer-reviewed
Theorem~\ref{thm:prafun}. The intrinsic identification $\pi_!A\cong B'$
(Proposition~\ref{prop:sigma}) is now proved for all $S,T$; all cardinality and equivariance
witnesses quoted above were machine-checked.

\begin{thebibliography}{9}
\bibitem{SGF}
D.~I.\ Spivak, R.\ Garner, and A.\ Fairbanks,
\emph{Functorial Aggregation},
arXiv:2111.10968, 2021.
\bibitem{NiuSpivak}
N.\ Niu and D.~I.\ Spivak,
\emph{Polynomial Functors: A Mathematical Theory of Interaction},
arXiv:2312.00990, 2023.
\end{thebibliography}

\end{document}
